\begin{pastexam}[label={px:ust-112-ee-ds-23}]{臺聯大電機類 112 學年度 第 23 題}
（問答；10 分）

Given the following pseudo code of Dijkstra algorithm, please point out the errors of the pseudo
code. (10 points)
\begin{lstlisting}[basicstyle=\ttfamily\small, frame=single, columns=fullflexible,
    keepspaces=true, mathescape=true, xleftmargin=3em, xrightmargin=3em]
1   procedure Dijkstra(Graph, source):
2       create vertex set Q
3       for each vertex v in Graph:
4           distance[v] $\leftarrow$ INFINITY
5           previous[v] $\leftarrow$ UNDEFINED
6           if source!=v:
7               add v to Q
8       distance[source] $\leftarrow$ 0
9       while Q is not empty:
10          u $\leftarrow$ vertex in Q with max_distance[u]
11          remove u from Q
12          for each neighbor v of u:
13              alternative $\leftarrow$ distance[u] - length(u, v)
14              if alternative < distance[v]:
15                  alternative $\leftarrow$ distance[v]
16                  u $\leftarrow$ previous[v]
17      return distance[], previous[]
\end{lstlisting}
\end{pastexam}

\begin{pastsoln}
共有五個錯誤，在第 6--7、10、13、15、16 行。 Dijkstra 演算法的正確流程（資結書第~25.2~節、
虛擬碼~25.3）：所有頂點放進 $Q$；每次從 $Q$ 取出 \textit{distance} 最\textbf{小}的頂點 $u$，對它的
每條出邊 $(u, v)$ 做鬆弛（relax）：若 $\mathit{distance}[u] + w(u, v) < \mathit{distance}[v]$，就
更新 $\mathit{distance}[v]$ 並記 $v$ 的前驅是 $u$（資結書虛擬碼~25.2）。 邊的權重必須非負。

\begin{enumerate}
\item \textbf{第 6--7 行}：起點 \textit{source} 沒有放進 $Q$。 這樣起點永遠不會被取出，它的出邊
不會被鬆弛；其他頂點的 \textit{distance} 一直是 INFINITY，算不出任何最短距離。 應該把每個頂點都
加進 $Q$：刪掉第 6 行的條件，第 7 行改成對每個 $v$ 都做 add $v$ to $Q$。
\item \textbf{第 10 行}：應該取 \textit{distance} \textbf{最小}的頂點，不是最大。 Dijkstra 的正確性
靠的就是「$Q$ 裡距離最小的頂點，它的距離已經是最短距離」（資結書定理~25.6）；取最大的就沒有這個
性質。 改成 u $\leftarrow$ vertex in $Q$ with min distance[u]。
\item \textbf{第 13 行}：經過 $u$ 到 $v$ 的路徑長是 $\mathit{distance}[u]$ \textbf{加上}邊長，不是減。
改成 alternative $\leftarrow$ distance[u] + length(u, v)。
\item \textbf{第 15 行}：方向寫反了。 找到更短的路徑時要更新的是 $\mathit{distance}[v]$，改成
distance[v] $\leftarrow$ alternative。 原本的寫法把 alternative 改掉，distance 陣列從頭到尾都不會變。
\item \textbf{第 16 行}：同樣寫反了。 要記錄的是「$v$ 的前一個頂點是 $u$」，改成
previous[v] $\leftarrow$ u。 原本的寫法還會在迴圈中途改掉 $u$，之後第 12 行的迴圈和第 13 行的
distance[u] 都用錯頂點。
\end{enumerate}

改正後：
\begin{lstlisting}[basicstyle=\ttfamily\small, frame=single, columns=fullflexible,
    keepspaces=true, mathescape=true, xleftmargin=3em, xrightmargin=3em]
1   procedure Dijkstra(Graph, source):
2       create vertex set Q
3       for each vertex v in Graph:
4           distance[v] $\leftarrow$ INFINITY
5           previous[v] $\leftarrow$ UNDEFINED
6           add v to Q
7       distance[source] $\leftarrow$ 0
8       while Q is not empty:
9           u $\leftarrow$ vertex in Q with min distance[u]
10          remove u from Q
11          for each neighbor v of u:
12              alternative $\leftarrow$ distance[u] + length(u, v)
13              if alternative < distance[v]:
14                  distance[v] $\leftarrow$ alternative
15                  previous[v] $\leftarrow$ u
16      return distance[], previous[]
\end{lstlisting}
第 14 行的比較（alternative $<$ distance[v]）本身沒有錯，保留。 $Q$ 用陣列時每次找最小值要
$O(|V|)$，總時間 $O(|V|^2)$；用二元堆積時是 $O((|V| + |E|) \log |V|)$。
\end{pastsoln}
